\chapter{Conclusiones}
\label{chapter:8-Conclusiones}

Se logró desarrollar satisfactoriamente un chatbot educativo basado en el contenido de OpenFing. Mediante nuestra prueba de concepto logramos mostrar que el alto volumen de datos, el gran cómputo requerido y el tiempo que demanda el preprocesamiento de las clases no son un impedimento. A pesar de las dificultades mencionadas, la solución es viable y funciona como es esperada.

Los resultados concretos obtenidos respaldan esta conclusión. En la evaluación de transcripción automática, el modelo Whisper \textit{medium} con especificación explícita del idioma español alcanzó una WER promedio inferior a 0.05 en todos los conjuntos de prueba evaluados, equivalente a menos de cinco errores por cada cien palabras. El pipeline de transcripción fue aplicado al corpus completo de OpenFing, que comprende 239 cursos y 4.339 clases, utilizando \textit{faster-whisper} con CTranslate2, lo que redujo los tiempos de procesamiento entre dos y cuatro veces respecto a la implementación original de Whisper. La generación de respuestas se implementó con el modelo \textit{meta-llama/Meta-Llama-3.1-8B-Instruct} (Llama~3.1~8B) accedido mediante la API de inferencia de HuggingFace, que demostró ser adecuado para el dominio académico en español sin requerir infraestructura de GPU propia.

Para el estudiante, la solución brinda respuestas mayoritariamente correctas con relación a las clases y cursos en los que se basa, en tiempo y forma razonable. Puede resultar muy útil a la hora de estudiar, ya que facilita el aprendizaje, agiliza la asimilación de conceptos y apoya el estudio autónomo, consiguiendo respuestas inmediatas sobre consultas específicas y complementando los recursos actuales.

Para la facultad, esta prueba de concepto representa una oportunidad para escalar el acceso al conocimiento. La implementación logró consolidar material educativo disperso en los múltiples videos en una interfaz conversacional única. Los recursos computacionales y de procesamiento necesarios para mantener el proyecto ya están disponibles en la infraestructura de la facultad, de modo que su operación continua no debería representar una carga significativa.

Respecto a los objetivos específicos planteados en la \autoref{section:1-Objetivos del Proyecto}, el proyecto los abordó en su totalidad aunque con distinta profundidad. El primer objetivo, diseñar e implementar un pipeline automatizado de extracción, descarga, transcripción e indexación, fue cumplido con un sistema modular de cinco componentes orquestados mediante Prefect. El segundo objetivo, evaluar herramientas ASR para el español académico rioplatense, fue satisfecho con la evaluación cuantitativa y cualitativa de los modelos Whisper \textit{base}, \textit{small} y \textit{medium}, concluyendo en la selección de \textit{medium} con especificación de idioma. El tercer objetivo, implementar un chatbot RAG con referencia al instante exacto del video, fue cumplido: el sistema recupera fragmentos de transcripción e incluye metadatos de temporización en cada respuesta. El cuarto objetivo, evaluar la calidad del pipeline RAG mediante experimentos de pregunta-respuesta, fue abordado parcialmente: la evaluación de transcripción quedó completa y documentada, mientras que la evaluación del componente conversacional sobre el corpus completo constituye la principal área pendiente. El quinto objetivo, desarrollar el sistema como software de código abierto con arquitectura extensible, fue logrado con una arquitectura de microservicios que permite incorporar nuevas bases de conocimiento sin rediseño estructural.

Si bien el prototipo actual cumple su función demostrativa, existen áreas claras de mejora: refinamiento de la precisión de respuestas, implementación de métricas de uso, integración con sistemas académicos existentes, y expansión del corpus con material de diversas asignaturas. Estas mejoras son técnicamente factibles y mantienen la filosofía de eficiencia de recursos del proyecto original.

En términos de autocrítica, el alcance del procesamiento resultó más amplio de lo inicialmente estimado. El preprocesamiento del corpus completo de OpenFing demandó un tiempo de cómputo significativo que condicionó los ciclos de experimentación con distintas configuraciones. Una estimación más precisa de los costos computacionales en la etapa de planificación habría permitido distribuir mejor el tiempo disponible. Por otra parte, la evaluación del componente de pregunta-respuesta no alcanzó la profundidad prevista, ya que el tiempo disponible priorizó la construcción y estabilización del pipeline completo sobre la fase experimental.

El proyecto muestra que las barreras para implementar herramientas educativas de tipo \acrshort{rag} basadas en \acrshort{ia} son principalmente metodológicas, no económicas, lo que abre posibilidades para democratizar el acceso a asistentes pedagógicos inteligentes en contextos académicos.

\section{Trabajo a futuro}
\label{chapter:7-Trabajo a futuro}
En esta sección se detallarán algunos de los avances necesarios para mejorar el proyecto actual. 

Para cumplir nuestra visión y que los estudiantes utilicen esta herramienta en el desarrollo de sus cursos, es necesario integrar el proyecto dentro de un entorno institucional, mantenido por o con el apoyo directo de la URI (Unidad de Recursos Informáticos) y OpenFing. La integración institucional aporta la sostenibilidad, estabilidad y el respaldo necesario a largo plazo. Durante el transcurso del proyecto, el sistema estuvo alojado en la plataforma "Mi Nube" de Antel como solución provisional mientras se avanza hacia la implementación de la infraestructura definitiva. Una posibilidad es evaluar si la URI pudiera asumir la responsabilidad del despliegue y mantenimiento del entorno operativo para este sistema. En \autoref{Estimaciones de lo que significaría para OpenFing} se presenta un documento entregado a la UEFI, donde se explican los recursos informáticos utilizados por el grupo y los procesos diseñados. Otro punto clave también sería un rediseño del actual procedimiento para subir material a OpenFing para que las clases nuevas se incorporen automáticamente al sistema. De esta manera se puede mantener el material actualizado a lo largo del tiempo sin necesidad de un trabajo manual. 

Ya cubierto el contenido en OpenFing, aún es posible aumentar el contenido para que otras facultades logren integrarse en el sistema. Las otras facultades de la UdelaR están comenzando a subir sus clases de manera online, desarrollando herramientas similares a OpenFing. Durante nuestra exposición en Ingeniería de Muestra (ver \autoref{chapter:A-0-Anexo I-Ingeniería de Muestra}), se nos comentó lo útil que puede ser este proyecto para otras facultades. Cubrir otras facultades requeriría un esfuerzo acotado. Se necesitaría extender el pipeline de procesamiento para incorporar clases de otras fuentes y añadir un filtrado por facultad de origen.

Este filtrado resultaría esencial en caso de añadir clases de distintas facultades debido a que muchos temas son abordados en cursos de diferentes facultades. Sin embargo, sí se realizan distinciones con respecto al enfoque, nivel de profundidad y contextos académicos. Segmentar la información por facultad permite que el chatbot acceda únicamente a las fuentes más pertinentes para cada caso. De esta forma, las respuestas entregadas no solo serán más precisas y relevantes, sino que también mantendrán coherencia con la perspectiva institucional y disciplinar correspondiente. Esto contribuye a una experiencia más clara, confiable y adaptada a las necesidades específicas de cada estudiante o usuario.
 
Otro trabajo a futuro es evaluar distintas técnicas de recuperación sobre el sistema. Una de ellas es el \textit{Re-Ranking} de los fragmentos recuperados. La arquitectura RAG depende en gran medida de la calidad y la relevancia del contexto que se entrega al modelo generador, y ese contexto proviene de una búsqueda por similitud en el espacio de \textit{embeddings}. Si bien esos \textit{embeddings} se entrenan justamente para que la cercanía entre vectores refleje cercanía semántica entre los textos, esa cercanía semántica general no garantiza que un fragmento sea el más pertinente para responder una consulta puntual. A esto se suma que, en un esquema de recuperación de una sola etapa, cada texto se codifica en un vector de forma independiente de la consulta, lo que limita la capacidad de capturar interacciones finas entre la pregunta y el pasaje. El \textit{Re-Ranking} introduce una segunda etapa en la que un modelo evalúa conjuntamente cada par consulta-fragmento y reordena los resultados, priorizando los que presentan mayor relevancia para la consulta concreta. El impacto de esta técnica debe medirse de forma empírica. Para ello se puede construir un conjunto de consultas con los fragmentos relevantes anotados, en la línea del conjunto de preguntas de evaluación que se elaboró para los experimentos del \autoref{chapter:6-Experimentación}, y comparar la relevancia del contexto recuperado con y sin reordenamiento. En \cite{irany2026hybrid}, los autores implementan una estrategia híbrida de recuperación y \textit{Re-Ranking} en un sistema RAG orientado al dominio biomédico, obteniendo mejoras en la calidad del contexto recuperado y resultados favorables en tareas de generación fundamentada en evidencia.

Ese mismo conjunto de consultas con fragmentos relevantes anotados habilitaría además un ajuste fino del modelo de \textit{embeddings} sobre el dominio de OpenFing. Los modelos de recuperación de propósito general se entrenan sobre corpus amplios y genéricos, por lo que su representación del vocabulario técnico de los cursos de ingeniería y del español académico rioplatense es necesariamente aproximada. Con un volumen suficiente de pares consulta-fragmento relevante, se podría adaptar el modelo para que acerque en el espacio vectorial las consultas de los estudiantes a los fragmentos que efectivamente las responden, lo que mejoraría la calidad de la recuperación en la primera etapa, antes de cualquier reordenamiento posterior.

En una línea complementaria, la recuperación actual se basa exclusivamente en similitud vectorial densa, lo que captura bien la similitud semántica entre la consulta y los fragmentos indexados, pero puede no rendir de forma óptima ante términos técnicos exactos, siglas o nombres de asignaturas poco frecuentes en el espacio vectorial. Una estrategia de búsqueda híbrida que combine la recuperación densa actual con un componente de recuperación léxica dispersa, como BM25, podría mejorar la precisión en este tipo de consultas. En ese esquema, los fragmentos recuperados por ambos métodos se fusionarían antes de la etapa de generación, aprovechando las fortalezas complementarias de cada enfoque.

Una mejora transversal que habilitaría varias de las anteriores es la incorporación de un mecanismo de retroalimentación explícita del usuario. Agregar controles simples como valoraciones de calidad sobre las respuestas del asistente permitiría acumular datos reales sobre qué respuestas son consideradas útiles y cuáles no. Este conjunto de interacciones etiquetadas no solo facilitaría el monitoreo continuo de la calidad del sistema en producción, sino que constituiría la materia prima necesaria para procesos de ajuste fino posteriores.

Otro trabajo a futuro relevante es el ajuste fino (\textit{fine-tuning}) del modelo de lenguaje sobre el dominio académico uruguayo. En una arquitectura RAG, la etapa de recuperación determina qué fragmentos del corpus llegan al modelo, pero es la etapa de generación la que decide cómo sintetizarlos en una respuesta. El modelo actual, Llama 3.1 8B, opera sin adaptación de dominio específica, por lo que puede producir respuestas que mezclan el contexto recuperado con conocimiento general del preentrenamiento, lo que aumenta el riesgo de inexactitudes en temas muy específicos de la FING. Ajustar el modelo sobre interacciones reales etiquetadas como correctas o incorrectas podría contribuir a que tienda a mantenerse más fiel al contexto recuperado, a reconocer con mayor frecuencia cuándo ese contexto es insuficiente para responder, y a emplear la terminología propia de los cursos de ingeniería en la UdelaR.

Técnicas de ajuste fino eficiente como LoRA \cite{hu2022lora} y QLoRA \cite{dettmers2023qlora} permiten adaptar modelos de la escala de Llama 3.1 8B con recursos de cómputo moderados, sin necesidad de reentrenar todos los parámetros del modelo. El proceso requeriría acumular primero datos de calidad a partir de las interacciones reales de los usuarios y luego aplicar el ajuste fino sobre los pares de consulta y respuesta validados. La mejora esperada no es en el recall de la recuperación vectorial, sino en la fidelidad y coherencia de las respuestas generadas a partir de los fragmentos ya recuperados.

Otra mejora a futuro consiste en incorporar \textit{streaming} de tokens en la generación de respuestas. Actualmente el sistema invoca al modelo de lenguaje de forma síncrona y solo envía la respuesta completa al frontend una vez finalizada la generación. En consultas con respuestas extensas, este enfoque puede aumentar la percepción de latencia por parte del usuario, que no recibe ninguna señal de progreso hasta que el proceso completo termina. LangChain soporta de forma nativa la emisión incremental de tokens, lo que permitiría adaptar tanto el backend como el frontend para mostrar la respuesta a medida que se genera, de forma similar a otras interfaces conversacionales basadas en LLM.

El corpus actual indexa únicamente el contenido textual de las transcripciones automáticas. Sin embargo, las clases de OpenFing incluyen diapositivas, diagramas y escritura en pizarra que contienen información complementaria no capturada en el texto transcripto. Extender el índice vectorial para incorporar representaciones de imágenes mediante modelos de visión y lenguaje podría enriquecer el contexto disponible para el asistente, en particular en consultas sobre conceptos que se explican principalmente a través de esquemas visuales o fórmulas escritas en pizarra.


